        \btabularx
        \hline
        \multicolumn{4}{|l|}{\bf value mode attributes} \\
        \hline
        domain   & 3PL     & clock   & resolve the clock domains of assigned values \\     
        readonly & 3PL     & boolean & set to read only \\     
%        maxdelay & netlist & log     & maximum delay    \\     
        \hline
        \btabularxend
        
        {\bf These attributes are case-insensitive.}

        {\bf domain} resolves the clock domains of any values assigned to the value
        variable to the specified clock domain. If assigned {\bf null} the current
        clock domain is selected. If the values are already resolved their clock
        domains are compared with the intended domain. If the assigned clock domain
        agrees with that within the values then there is no further action,
        otherwise a fatal error occurs. If a value has not been assigned, the clock
        domain will be set and a future assignment must have the same clock domain.
        The attribute will affect all variables contained in the value expression.
        Setting the clock domain has a particular use for value variables which are
        outputs of an {\bf import()} \autorefph{sec:ipimport} procedure. Note that if
        this attribute is read back a null is returned.
        
        {\bf readonly} true specifies that from now on this variable cannot
        be assigned to. This attribute cannot be assigned false.
        
%        {\bf maxdelay} is a place-and-route constraint. This places an
%        upper limit on the propagation time through the signal net or nets
%        comprising this value.
